% ============================================================================
%  PHẦN ĐẦU: Lời cảm ơn · Tóm tắt · Mục lục · Danh sách hình/bảng · Từ viết tắt
% ============================================================================

\section*{Lời cảm ơn}
\addcontentsline{toc}{section}{Lời cảm ơn}

Nhóm 7 xin gửi lời cảm ơn chân thành đến ThS. Nguyễn Hữu Quyền, giảng viên môn
\emph{Lập trình Python cho Máy học} (CS116), đã định hướng đề tài, góp ý về cách
thiết kế thí nghiệm và nhấn mạnh yêu cầu \emph{nêu rõ giới hạn dữ liệu, không khẳng
định vượt quá phạm vi thực nghiệm} --- yêu cầu đã định hình cách nhóm xây dựng và
trình bày toàn bộ đồ án này.

Nhóm cũng cảm ơn Phòng thí nghiệm NLP@UIT đã công bố bộ dữ liệu UIT-ViQuAD và mô hình
ViSoBERT, cùng cộng đồng mã nguồn mở HuggingFace, PyTorch và scikit-learn đã cung cấp
hạ tầng để một nhóm sinh viên có thể huấn luyện và đánh giá mô hình Transformer trên
máy tính cá nhân.

\vspace{1em}
\hfill TP. Hồ Chí Minh, tháng 9 năm 2026

\hfill Nhóm 7

\clearpage
% ----------------------------------------------------------------------------
\section*{Tóm tắt}
\addcontentsline{toc}{section}{Tóm tắt}

Đồ án xây dựng và đánh giá một hệ thống \textbf{đọc hiểu máy trích xuất} (extractive
Machine Reading Comprehension --- MRC) cho tiếng Việt: cho một đoạn văn và một câu hỏi,
hệ thống chỉ ra đoạn con của văn bản chứa câu trả lời, hoặc kết luận rằng đoạn văn
\emph{không} chứa câu trả lời. Dữ liệu là \textbf{UIT-ViQuAD~2.0} (28.454 câu hỏi huấn
luyện, trong đó 32,39\% là câu không có đáp án).

Nhóm so sánh bốn mô hình, mỗi mô hình cô lập một biến thực nghiệm: (1) baseline truy
hồi câu bằng \textbf{TF-IDF}; (2) \textbf{XLM-R} đã huấn luyện trên SQuAD~2.0 tiếng
Anh, dùng \emph{zero-shot}; (3) \textbf{mBERT} và (4) \textbf{ViSoBERT} được
\emph{fine-tune} trên toàn bộ tập train bằng GPU Apple Silicon (MPS). Toàn bộ pipeline
--- chuẩn hoá, metric, cửa sổ trượt, gán nhãn huấn luyện, đánh giá --- được tự cài
đặt theo phương pháp phát triển hướng kiểm thử (TDD) với 474 test.

\ifhasFullEval
Trên toàn bộ \valN{} câu của tập validation (checkpoint chọn trên dev khi có, xem
Mục~\ref{sec:experiments}), \textbf{mBERT fine-tune đạt EM \mbertVoneEM{} và F1
\mbertVoneFone{}}, vượt XLM-R zero-shot (EM \xlmrEM{} / F1 \xlmrFone{}), ViSoBERT (EM
\visoVoneEM{} / F1 \visoVoneFone{}) và TF-IDF (EM \baselineEM{} / F1 \baselineFone{}).
Phân tích tách nhóm cho thấy ba điều mà con số tổng che khuất: (i) khoảng cách EM--F1 của
baseline chứng minh hai metric đo hai kỹ năng khác nhau; (ii) mBERT hơn XLM-R ở hai kỹ năng
tách biệt: \emph{biết từ chối} (EM trên câu không có đáp án \mbertVoneNoAnsEM{} so với
\xlmrNoAnsEM{}) và \emph{xác định đúng biên đáp án} (EM trên câu có đáp án
\mbertVoneHasAnsEM{} so với \xlmrHasAnsEM{}), trong khi F1 trên câu có đáp án của hai mô
hình tương đương (\mbertVoneHasAnsFone{} so với \xlmrHasAnsFone{}) --- XLM-R zero-shot tìm
đúng vùng đáp án nhưng chưa học quy ước ranh giới đáp án của bộ dữ liệu; (iii) ViSoBERT ---
được pretrain trên văn bản mạng xã hội --- suy sụp về hành vi ``gần như luôn trả lời rỗng'',
điểm EM tổng bằng đúng tỉ lệ câu không có đáp án. Nhóm loại trừ lần lượt các nguyên nhân
cấu hình trước khi kết luận đây là giới hạn của mô hình trong thiết lập này, và gắn nó với
việc từ vựng \tokVocabViso{} token phân mảnh văn bản Wikipedia tiếng Việt nhiều hơn khoảng
1,6 lần so với mBERT. Việc tách EM/F1 theo câu có đáp án và câu không có đáp án theo đúng
quy ước đã có sẵn trong công cụ chấm chính thức của SQuAD~2.0; đóng góp của nhóm là áp dụng
quy ước đó một cách nhất quán ở mọi bảng, không phải một phép đo mới.
\else
Trên mẫu ngẫu nhiên 500 câu của tập validation (seed 42), \textbf{mBERT fine-tune đạt
EM \histMbertEM{} và F1 \histMbertFone{}}, vượt XLM-R zero-shot (EM \histXlmrEM{} / F1
\histXlmrFone{}), ViSoBERT (EM \histVisoEM{} / F1 \histVisoFone{}) và TF-IDF (EM
\histBaselineEM{} / F1 \histBaselineFone{}). Phân tích tách nhóm cho thấy ba điều mà con số
tổng che khuất: (i) khoảng cách EM--F1 của baseline chứng minh hai metric đo hai kỹ
năng khác nhau; (ii) mBERT hơn XLM-R ở hai kỹ năng tách biệt: \emph{biết từ chối} (EM trên
câu không có đáp án \histMbertNoAnsEM{} so với \histXlmrNoAnsEM{}) và \emph{xác định đúng
biên đáp án} (EM trên câu có đáp án \histMbertHasAnsEM{} so với \histXlmrHasAnsEM{}), trong
khi F1 trên câu có đáp án của hai mô hình tương đương (\histMbertHasAnsFone{} so với
\histXlmrHasAnsFone{}) --- XLM-R zero-shot tìm đúng vùng đáp án nhưng chưa học quy ước ranh
giới đáp án của bộ dữ liệu; (iii) ViSoBERT --- được pretrain trên văn
bản mạng xã hội --- suy sụp về hành vi ``gần như luôn trả lời rỗng'', điểm EM tổng bằng
đúng tỉ lệ câu không có đáp án của mẫu. Nhóm loại trừ lần lượt các nguyên nhân cấu hình
trước khi kết luận đây là giới hạn của mô hình trong thiết lập này, và gắn nó với việc
từ vựng \tokVocabViso{} token phân mảnh văn bản Wikipedia tiếng Việt nhiều hơn khoảng 1,6
lần so với mBERT. Việc tách EM/F1 theo câu có đáp án và câu không có đáp án theo đúng quy
ước đã có sẵn trong công cụ chấm chính thức của SQuAD~2.0; đóng góp của nhóm là áp dụng quy
ước đó một cách nhất quán ở mọi bảng, không phải một phép đo mới.
\fi

Mọi con số trong báo cáo được sinh từ một lần chạy thật và truy vết được về một tệp
\code{results/*.json} kèm commit, thời điểm và thiết bị. Các giả thuyết định lượng được
ghi \emph{trước} khi chạy và được đối chiếu công khai, kể cả những giả thuyết sai.
Hệ thống được đóng gói thành ứng dụng web Streamlit cho phép thử trực tiếp.

\vspace{0.6em}
\noindent\textbf{Từ khoá:} đọc hiểu máy, hỏi đáp trích xuất, tiếng Việt, UIT-ViQuAD~2.0,
Transformer, mBERT, XLM-R, ViSoBERT, câu hỏi không có đáp án, tái lập thực nghiệm.

\clearpage
% ----------------------------------------------------------------------------
{\setstretch{1.1}
\tableofcontents
\clearpage
\listoffigures
\addcontentsline{toc}{section}{Danh sách hình vẽ}
\vspace{1.5em}
\listoftables
\addcontentsline{toc}{section}{Danh sách bảng}
}

\vspace{1.5em}
% ----------------------------------------------------------------------------
\section*{Danh mục từ viết tắt}
\addcontentsline{toc}{section}{Danh mục từ viết tắt}

\begin{center}
\renewcommand{\arraystretch}{1.15}
\begin{tabularx}{\textwidth}{@{}l L@{}}
\toprule
\textbf{Viết tắt} & \textbf{Ý nghĩa} \\
\midrule
MRC       & Machine Reading Comprehension --- đọc hiểu máy \\
QA        & Question Answering --- hỏi đáp \\
EM        & Exact Match --- khớp chính xác \\
F1        & Trung bình điều hoà của precision và recall (ở đây tính trên token) \\
TF-IDF    & Term Frequency -- Inverse Document Frequency \\
BERT      & Bidirectional Encoder Representations from Transformers \\
mBERT     & Multilingual BERT (\code{bert-base-multilingual-cased}) \\
XLM-R     & XLM-RoBERTa --- mô hình đa ngữ của Conneau và cộng sự \\
ViSoBERT  & Vietnamese Social media BERT (\code{uitnlp/visobert}) \\
SQuAD     & Stanford Question Answering Dataset \\
UIT-ViQuAD & UIT Vietnamese Question Answering Dataset \\
MPS       & Metal Performance Shaders --- backend GPU của PyTorch trên Apple Silicon \\
TDD       & Test-Driven Development --- phát triển hướng kiểm thử \\
CI        & Confidence Interval --- khoảng tin cậy \\
CLS       & Token đặc biệt ở đầu chuỗi đầu vào của BERT, dùng để biểu diễn ``không có đáp án'' \\
\bottomrule
\end{tabularx}
\end{center}
